% Curriculum vitae — Elie Attal
% Compile with: pdflatex cv.tex   (or: latexmk -pdf cv.tex)

\documentclass[11pt,a4paper]{article}

% ---------- Packages ----------
\usepackage[utf8]{inputenc}
\usepackage[T1]{fontenc}
\usepackage{lmodern}
\usepackage{microtype}
\usepackage[margin=2.2cm]{geometry}
\usepackage{amsmath,amssymb}
\usepackage{enumitem}
\usepackage{titlesec}
\usepackage{needspace}
\usepackage{xcolor}
\usepackage[hidelinks]{hyperref}

\pagestyle{empty}
\setlength{\parindent}{0pt}
\setlength{\parskip}{0pt}

% ---------- Section style ----------
\definecolor{rule}{gray}{0.6}
\titleformat{\section}{\large\scshape}{}{0pt}{}[{\color{rule}\titlerule}]
\titlespacing*{\section}{0pt}{14pt}{8pt}
% never leave a heading alone at the bottom of a page
\let\oldsection\section
\renewcommand{\section}{\Needspace{5\baselineskip}\oldsection}

% ---------- Helpers ----------
% \entry{left column (dates)}{right column (content)}
\newlength{\leftcol}
\setlength{\leftcol}{2.8cm}
\newcommand{\entry}[2]{%
  \makebox[\leftcol][l]{\small #1}%
  \parbox[t]{\dimexpr\linewidth-\leftcol\relax}{#2}\par
  \vspace{5pt}}

% \pub{title}{coauthors}{venue}{arXiv id}
\newcommand{\pub}[4]{%
  \item \textbf{#1}\\
  {\small #2. #3.
  \href{https://arxiv.org/abs/#4}{arXiv:#4}}}

\newcommand{\invited}{\textit{(invited)}}
\newcommand{\prize}[1]{\textit{#1}}

% =====================================================================
\begin{document}

% ---------- Header ----------
{\LARGE\bfseries Elie Attal}\par
\vspace{4pt}
{\small
  \href{mailto:elie.attal@polytechnique.edu}{elie.attal@polytechnique.edu}
  \;$\cdot$\;
  +33 6 25 70 15 51
  \;$\cdot$\;
  \href{https://eattal.github.io}{eattal.github.io}
  \;$\cdot$\;
  \href{https://arxiv.org/search/math?searchtype=author&query=Attal,+E}{arXiv}
  \;$\cdot$\;
  \href{https://scholar.google.com/citations?user=jp5K33sAAAAJ}{Google Scholar}
}

% ---------- Education ----------
\section{Education}
\entry{2024 -- present}{%
  \textbf{PhD in applied mathematics}\\ 
  \textit{CMAP, École Polytechnique}\\
  Advisors: E. Abi Jaber and M. Rosenbaum.\\
  Thesis: Interacting agents, jump processes and Volterra processes.}
\entry{2023 -- 2024}{%
  \textbf{M2 Probability and Finance}\\
   \textit{Sorbonne Université / École Polytechnique}}
\entry{2020 -- 2023}{%
  \textbf{Ingénieur polytechnicien master's degree}\\
  \textit{École Polytechnique}\\
  Master's degree in applied mathematics.}

% ---------- Professional experience ----------
\section{Professional experience}
\entry{4 months, 2026}{\textbf{Quantitative research intern}\\
\textit{Bloomberg L.P., New York}\\
  {}Developed numerical methods for path-dependent option pricing under rough volatility models. Carried out as part of the Bloomberg Quantitative Finance PhD Fellowship.}
\entry{6 months, 2024}{\textbf{Quantitative research intern}\\
\textit{Qube Research \& Technologies, Paris}\\
  {}Studied the limiting eigenvector overlaps between large random matrices and their submatrices. Published as two papers in \textit{Journal of Physics A: Mathematical and Theoretical}.}
\entry{5 months, 2023}{\textbf{Quantitative research intern}\\
\textit{DRW Holdings, London}\\
  {}Designed a  quoting model for RFQ platforms combining machine learning and probability.}
\entry{3 months, 2022}{\textbf{Quantitative research intern}\\
\textit{ITO33, Paris}\\
{}Extended the Black--Scholes framework to price convertible bonds under takeover risk. Implemented the model in C++ in the company's pricing software.}

% ---------- Awards ----------
\section{Awards}
\entry{2026}{\textbf{QRT PhD Prize}\\
  For the talk \emph{Affine structure of the Brownian signature} at the London--Paris Bachelier Workshop on Mathematical Finance.}
\entry{2025}{\textbf{Bloomberg Quantitative Finance PhD Fellowship}}

% ---------- Publications ----------
\section{Publications}
\begin{itemize}[leftmargin=*,itemsep=4pt]
  \pub{Simulating integrated Volterra square-root processes and Volterra Heston models via inverse Gaussian}
      {with E. Abi Jaber}
      {\emph{Finance and Stochastics}, to appear, 2026}
      {2504.19885}
  \pub{From hyper roughness to jumps as $H \to -1/2$}
      {with E. Abi Jaber and M. Rosenbaum}
      {\emph{The Annals of Applied Probability} 36(4), 3635--3660, 2026}
      {2503.16985}
  \pub{Eigenvector overlaps of random covariance matrices and their submatrices}
      {with R. Allez}
      {\emph{Journal of Physics A: Mathematical and Theoretical} 58, 205003, 2025}
      {2501.08768}
  \pub{Interlacing eigenvectors of large Gaussian matrices}
      {with R. Allez}
      {\emph{Journal of Physics A: Mathematical and Theoretical} 58, 045003, 2025}
      {2409.17086}
\end{itemize}

% ---------- Preprints ----------
\section{Preprints}
\begin{itemize}[leftmargin=*,itemsep=4pt]
  \pub{Stochastic control with signatures via Riccati equations on the tensor algebra}
      {with E. Abi Jaber and D. Sotnikov}
      {2026}
      {2607.03986}
  \pub{Affine structure of the Brownian signature}
      {with E. Abi Jaber and D. Sotnikov}
      {2026}
      {2606.29622}
  \pub{Volterra clocks and their pure-jump limits: hitting times of curved boundaries}
      {with E. Abi Jaber and A. Søjmark}
      {2026}
      {2605.29073}
  \pub{Efficient simulation of Hawkes processes using their affine Volterra structure}
      {with E. Abi Jaber and D. Sotnikov}
      {2025}
      {2511.13554}
\end{itemize}

% ---------- Talks ----------
\section{Talks}
\textbf{Conferences and workshops}\par\vspace{4pt}
\entry{Sept.\ 2026}{London--Paris Bachelier Workshop on Mathematical Finance, Paris \invited\\
  \emph{Affine structure of the Brownian signature} --- \prize{QRT PhD Prize}}
\entry{Sept.\ 2026}{International Conference on Computational Finance, Oxford\\
  \emph{Simulating integrated Volterra square-root processes and Volterra Heston models via inverse Gaussian}}
\entry{Apr.\ 2026}{Stochastic Control and Games for Risk \& Regulation, Hammamet\\
  \emph{Simulating affine Volterra processes via inverse Gaussian}}
\entry{Jan.\ 2026}{18th Bachelier Colloquium on Financial Mathematics and Stochastic Calculus, Métabief\\
  \emph{From hyper-roughness to jumps as $H \to -1/2$: the affine case and beyond}}
\entry{Sept.\ 2025}{Stochastics \& Computational Finance, Lisbon\\
  \emph{Simulating affine Volterra processes via inverse Gaussian}}
\entry{July 2025}{Vienna Congress on Mathematical Finance, Vienna\\
  \emph{From hyper roughness to jumps as $H$ goes to $-1/2$}}

\vspace{4pt}
\textbf{Seminars}\par\vspace{4pt}
\entry{Oct.\ 2026}{Seminar of the CMAP MathsFi group, École Polytechnique\\
  \emph{Volterra clocks: going beyond the affine structure}}
\entry{Jan.\ 2026}{CMAP PhD Seminar, École Polytechnique\\
  \emph{Introduction to non-Markovian diffusions in Mathematical Finance}}
\entry{Dec.\ 2025}{Seminar of the CMAP MathsFi group, École Polytechnique\\
  \emph{Simulating affine Volterra processes via inverse Gaussian: the iVi schemes}}
\entry{May 2025}{Local Seminar, Vega Institute, Moscow (online)\\
  \emph{From hyper roughness to jumps as $H \to -1/2$}}
\entry{Mar.\ 2025}{Bachelier Seminar (PhD session), IHP, Paris\\
  \emph{From hyper roughness to jumps as $H \to -1/2$}}
\entry{Mar.\ 2025}{Seminar of the CMAP MathsFi group, École Polytechnique\\
  \emph{From hyper roughness to jumps as $H \to -1/2$}}

% ---------- Teaching ----------
\section{Teaching}
\entry{2026 -- 2027}{\textbf{École Polytechnique}\\
  Supervision of M1 group projects: \emph{Stability of eigenvectors of covariance matrices in finance};
  \emph{Heat equation and Cole--Hopf transform: entire series representations and numerical schemes}}
\entry{2025 -- 2026}{\textbf{École Polytechnique}\\
  Markov processes and applications (M1), tutorial classes\\
  Probability (L3), tutoring\\[3pt]
  \textbf{ENSAE}\\
  Supervision of an M1 group project: \emph{Toward a new stochastic model for integrated variance:
  insights from SPX and VIX}}
\entry{2024 -- 2025}{\textbf{École Polytechnique}\\
  Markov processes and applications (M1), tutorial classes\\
  Probability (L3), tutoring\\
  Supervision of an M1 internship: \emph{Numerical schemes for Riccati--Volterra equations}}

% ---------- Skills ----------
\section{Skills}
\entry{Languages}{French (native), English (fluent), Spanish (conversational)}
\entry{Programming}{Python (numpy, scipy, pandas, scikit-learn), SQL, Git, LaTeX, Power BI}

\end{document}
